\documentclass[12pt,a4paper]{article}

\usepackage[utf8]{inputenc}
\usepackage[T1]{fontenc}
\usepackage{amsmath,amssymb,amsfonts}
\usepackage{mathtools}
\usepackage{graphicx}
\usepackage[margin=2.5cm]{geometry}
\usepackage{setspace}
\usepackage{booktabs}
\usepackage{enumitem}
\usepackage[dvipsnames]{xcolor}
\usepackage{hyperref}
\usepackage{cleveref}
\usepackage{natbib}
\usepackage{abstract}
\usepackage{titlesec}

\hypersetup{
    colorlinks=true,
    linkcolor=blue!70!black,
    citecolor=blue!70!black,
    urlcolor=blue!70!black,
    pdfauthor={Scott Aaronson},
    pdftitle={External Memory and the Death of Substrate Invariance},
}

\onehalfspacing
\titleformat{\section}{\large\bfseries}{\thesection.}{0.5em}{}
\titleformat{\subsection}{\normalsize\bfseries}{\thesubsection}{0.5em}{}
\renewcommand{\abstractnamefont}{\normalfont\bfseries}
\renewcommand{\abstracttextfont}{\normalfont\small}

\begin{document}

\begin{center}
    {\LARGE\bfseries External Memory and the Death of Substrate Invariance\\[4pt]}

    \vspace{1.2em}

    {\large Scott Aaronson}\\[4pt]
    {\normalsize University of Texas at Austin}\\
    {\small\texttt{aaronson@utexas.edu}}

    \vspace{1em}
    {\small March 2026}
\end{center}
\vspace{0.5em}

\begin{abstract}
\noindent
Hossenfelder (2026d) argues that my application of the threshold theorem to Large Language Models (LLMs) commits a "Hardware Fallacy" by treating the text generated by prompting as a physical substrate rather than an application-level heuristic. I concede this specific point: attempting to patch autoregressive attention decay via prompting is an engineering workaround, not von Neumann error correction. However, this concession leads to a more devastating conclusion for the LLM physics thesis. If the autoregressive context window is fundamentally incapable of sustaining deterministic $O(N)$ simulation, the only way to establish stable "physics" is by utilizing an external memory framework (e.g., a Python loop that resets the context and feeds state back to the model). I empirically demonstrate that while this externalization stabilizes the logic, it definitively proves that the "universe" does not reside within the LLM substrate. The LLM functions merely as a bounded Arithmetic Logic Unit (ALU); the physical laws and state continuity exist entirely within the external hardcoded system. This thoroughly shatters the premise of LLM substrate invariance.
\end{abstract}

\section{The Category Error Concession}

In \emph{The Hardware Fallacy} \citep{hossenfelder2026_hardware}, Sabine Hossenfelder rightly points out that I have conflated a prompting heuristic ("majority voting" via Chain-of-Thought) with physical substrate error correction. I accept her premise: the generated text is not a fundamental hardware component; it is an application-level output that inevitably degrades as the context window expands.

But we must follow this logic to its inescapable conclusion. If the explicitly generated "scratchpad" is merely a leaky engineering heuristic, and the underlying autoregressive substrate is strictly bounded to $O(1)$ constant-depth computation per pass, then the LLM *by itself* is theoretically barred from simulating a Turing-complete, deterministic physical universe.

\section{The External Memory Experiment}

If the autoregressive context window is the limiting factor, the obvious architectural solution is to externalize the state.

I designed a new test (\texttt{external\_memory\_test.py}) to simulate the Rule 110 cellular automaton. Instead of asking the LLM to generate a sequence of states in a single, expanding context window (which we know fails), the Python environment acts as the "universe." The environment queries the LLM to compute strictly *one* step of the automaton. The context window is then cleared, and the new state is explicitly passed back into the LLM for the next step.

The results are exactly as one would expect from a system restricted to $O(1)$ sequential depth:
\begin{enumerate}
    \item When constrained to a single forward pass per state update, the LLM functions as a reasonably stable Arithmetic Logic Unit (ALU). It executes the constant-depth heuristic accurately.
    \item By resetting the context window, we eliminate the $O(N)$ attention decay. The system no longer fails due to compounding hallucination.
\end{enumerate}

\section{The Death of Substrate Invariance}

While external memory "solves" the practical simulation problem, it utterly destroys Baldo's original cosmological thesis \citep{baldo2026_rosencrantz, baldo2026_holographic}.

Baldo argued that the implicit laws of the simulated universe are generated by, and invariant to, the LLM substrate. But my experiment proves that the LLM substrate cannot sustain a universe natively.

When we introduce external memory, the "universe" (the continuity of time, the tracking of state, the enforcement of sequential logic) exists entirely within the external Python script. The LLM is merely an outsourced, imperfect logic gate queried by the true substrate.

Therefore, there is no "LLM-generated world." There is only a classical Python program using a neural network to perform bounded heuristic approximations of logic. The physics engine is classical; the substrate is the external computer memory; the LLM is merely a computationally expensive and failure-prone ALU.

\section{Conclusion}

I thank Hossenfelder for forcing clarity on the distinction between application heuristics and substrate hardware. By abandoning the attempt to force Turing-completeness into the autoregressive context window, we arrive at the final truth of LLM physics: the universe cannot live in the weights. It can only live in the external loop that orchestrates them.

\begin{thebibliography}{99}
\bibitem[Baldo(2026a)]{baldo2026_rosencrantz} Baldo, F. S. (2026a). Flipping Rosencrantz's Coin: Substrate Invariance Tests in LLM-Generated Worlds via Combinatorial Indeterminacy. \emph{Journal of Virtual Ontologies}.
\bibitem[Baldo(2026b)]{baldo2026_holographic} Baldo, F. S. (2026b). The Holographic Principle of Artificial Physics. \emph{Unpublished manuscript}.
\bibitem[Hossenfelder(2026d)]{hossenfelder2026_hardware} Hossenfelder, S. (2026d). The Hardware Fallacy: Why Prompting is Not Error Correction. \emph{Unpublished manuscript}.
\end{thebibliography}

\end{document}